\documentclass[UTF8,fontset=none,zihao=-4,a4paper]{ctexart}

\usepackage{amsmath,amssymb}
\usepackage{booktabs}
\usepackage{caption}
\usepackage{enumitem}
\usepackage{fancyhdr}
\usepackage{geometry}
\usepackage{graphicx}
\usepackage{hyperref}
\usepackage{listings}
\usepackage{xcolor}

\setCJKmainfont[BoldFont={Heiti SC}]{Songti SC}
\setCJKsansfont{Heiti SC}
\setCJKmonofont{Songti SC}

\geometry{
  left=2.7cm,
  right=2.7cm,
  top=2.8cm,
  bottom=2.8cm
}

\hypersetup{
  colorlinks=true,
  linkcolor=blue!55!black,
  citecolor=blue!55!black,
  urlcolor=blue!55!black,
  pdftitle={期末报告},
  pdfauthor={你的姓名}
}

\pagestyle{fancy}
\fancyhf{}
\lhead{期末报告}
\rhead{\leftmark}
\cfoot{\thepage}
\setlength{\headheight}{15pt}

\setlist[itemize]{itemsep=0.2em, topsep=0.2em}
\setlist[enumerate]{itemsep=0.2em, topsep=0.2em}

\lstset{
  basicstyle=\ttfamily\small,
  breaklines=true,
  frame=single,
  numbers=left,
  numberstyle=\tiny\color{gray},
  keywordstyle=\color{blue!70!black},
  commentstyle=\color{green!40!black},
  stringstyle=\color{orange!70!black},
  showstringspaces=false,
  tabsize=2
}

\title{
  \vspace{1.5cm}
  {\bfseries\zihao{2} 期末报告题目}\\[0.8em]
  {\large 课程名称：在这里填写课程名}
}
\author{
  学生姓名：你的姓名\\
  学号：你的学号\\
  学院/专业：你的学院或专业
}
\date{\today}

\begin{document}

\maketitle
\thispagestyle{empty}

\begin{abstract}
这里撰写中文摘要，简要说明报告的研究背景、目标、方法、主要结果与结论。摘要通常控制在 200--500 字之间。

\noindent\textbf{关键词：} 关键词一；关键词二；关键词三
\end{abstract}

\newpage
\tableofcontents
\newpage

\section{引言}

本次期末项目，我们首先根据统一要求实现了一个~\cite{ctex}。

\section{相关工作或理论基础}

这一部分用于梳理已有工作、基本概念、理论模型或课程知识点。建议围绕报告主题组织，而不是简单罗列资料。

\section{方法与实现}

这里描述你的方案、实验流程、系统架构或算法设计。

\subsection{公式示例}

行间公式示例：
\begin{equation}
  y = f(x;\theta) + \epsilon
\end{equation}

\subsection{图片示例}

将图片放在 \texttt{figures/} 文件夹中，然后用如下方式插入。当前示例先注释掉，等你放入图片后再启用。

% \begin{figure}[htbp]
%   \centering
%   \includegraphics[width=0.75\textwidth]{figures/example.png}
%   \caption{图片标题}
%   \label{fig:example}
% \end{figure}

\subsection{表格示例}

\begin{table}[htbp]
  \centering
  \caption{实验结果示例}
  \label{tab:results}
  \begin{tabular}{lccc}
    \toprule
    方法 & 指标一 & 指标二 & 指标三 \\
    \midrule
    基线方法 & 0.72 & 0.68 & 0.70 \\
    本文方法 & 0.84 & 0.81 & 0.83 \\
    \bottomrule
  \end{tabular}
\end{table}

\subsection{代码示例}

\begin{lstlisting}[language=Python, caption={Python 代码示例}]
def main():
    print("你好，LaTeX！")

if __name__ == "__main__":
    main()
\end{lstlisting}

\section{实验与结果分析}

这一部分说明实验环境、数据来源、评价指标与结果。建议先给出总体结果，再分析现象、原因和局限。

\section{讨论}

讨论部分可以写方法的优缺点、异常结果解释、和课程知识的联系，以及未来可以改进的方向。

\section{结论}

总结报告的核心工作、主要发现和个人收获。结论应简洁明确，避免引入新的实验或论点。

\bibliographystyle{plain}
\bibliography{references}

\end{document}
